% 第5章 仿真：策略训练与可观测性验证

\chapter{仿真：策略训练与可观测性验证}\label{ch:sim_chapter}

\section{章节引言：仿真侧范式与本章组织}\label{sec:sim_intro}

\cref{ch:task_policy} 给出了任务策略的方法设计与共享网络架构（25 维观测、随机偏差训练、DINOv3-S 编码器加三个网络头），但留下了"具体怎么训"这一执行层问题。本章承接\cref{ch:task_policy} 的方法层，给出本文\textbf{仿真侧}的两条主线及其实验闭环：

\begin{itemize}
  \item \emph{任务策略仿真训练}（\cref{sec:sim_task_train}）：采用 HIL-SERL（RLPD $+$ SAC $+$ 人在回路干预）训练任务策略，配套 Actor-Learner 分布式架构。这一训练用于验证\cref{ch:problem_formulation} 提出的可观测性命题——由\cref{sec:sim_experiments} 中的 H1（偏差退化基线）、H2/H3（固定/随机偏差对照）、观测维度消融三组实验构成端到端闭环。\emph{消融实验在 PickCube 24D 上进行}以隔离编码器偏差因素;\cref{sec:obs_design} 描述的部署 $25$D 是 24D $+$ hand 通道的等价扩展，hand 通道不参与可辨识性论证、仅供运行时与备份策略的距离触发切换。
  \item \emph{备份策略仿真训练}（\cref{sec:backup_sim_train}）：采用纯在线 SAC $+$ 域随机化训练备份避让策略，配套 V1$\to$V2$\to$V3$\to$V3b 的奖励与几何演化。这条主线服务于\cref{ch:real_chapter} 的真机部署——backup 在真机上不重训，仿真 ckpt 直接迁移。
\end{itemize}

两条主线在仿真层面相互独立：任务策略关注"如何在编码器偏差下完成抓取"，备份策略关注"如何在人手接近时安全闪避"；运行时由距离判据硬切换协调，详见\cref{ch:real_chapter} 真机部署。\cref{sec:sim_experiments} 给出两条主线对应的全部仿真实验数据；\cref{sec:sim_to_real_bridge} 是本章末尾的简短桥接，把仿真已验证的"成功信号路径"过渡到真机部署。

\section{任务策略仿真训练（HIL-SERL）}\label{sec:sim_task_train}

\cref{sec:task_method_motivation} 的设计选择三与\cref{sec:il_hil} 已论证仿真侧采用 HIL-SERL 验证可观测性命题与随机偏差鲁棒训练的依据（高吞吐 Actor-Learner、HIL 干预补稀疏奖励、与真机栈架构同源）。本节给出实现框架：以 Luo 等人提出的 HIL-SERL~\cite{luo2025hilserl} 为主干，以人在回路干预为持续纠错信号，并通过 Actor-Learner 分布式架构支撑高吞吐异步采样。本节描述的框架\emph{在仿真闭环中收敛稳定}，配套的 H1–H3 实证见\cref{sec:sim_experiments};真机失败的方法论复盘见\cref{subsec:hilserl_failure}。

\paragraph{算法栈。}框架的底层是连续动作空间上的 Soft Actor-Critic~\cite{haarnoja2018sac};其上，Ball 等人提出的 RLPD~\cite{ball2023rlpd} 通过两个回放缓冲区的 50/50 混合采样为稀疏奖励任务提供 demo 引导——\emph{离线演示缓冲区}预先装入 50 条人类遥操作的成功 episode，\emph{在线缓冲区}收集策略与环境的实时交互，每一步梯度更新都从两侧各采一半，不修改 SAC 损失;本文沿用 RLPD 的高 UTD 比例与 LayerNorm 设计以稳定 critic 训练。HIL-SERL 则在 RLPD 之上进一步将人在回路干预作为训练数据的一等来源贯穿整个在线训练过程（干预数据的具体处理见下段）。

\paragraph{人在回路干预。}在线采样过程中，训练者可通过遥操作设备（键盘、SpaceMouse 等）介入策略正在执行的 episode 并直接给出期望动作;这一"在策略失败的区域由人类示教补充纠错"的思路可追溯到 Ross 等人提出的 DAgger~\cite{ross2011dagger}。干预期间产生的转移被统一追加到离线演示缓冲区。本文最初曾尝试仅保留"最终成功"的干预 episode，但在实现与训练过程中发现两点问题：(i)~转移以步为单位写入缓冲区而非按 episode 组织，按成功与否回溯分割需引入额外的边界判定逻辑，且容易因失败率较高而整体丢弃大量有效监督;(ii)~即便 episode 最终未能完成抓取，人类干预所给出的动作方向本身也是可供策略学习的纠错梯度信号。因此最终统一保留所有干预转移，不再依赖最终成功与否进行过滤。

\paragraph{分布式训练架构。}\cref{fig:hilserl_arch} 示意训练系统的两个异步进程。Actor 进程持有策略权重副本，负责与仿真环境交互、接收 HIL 干预、把转移写入在线缓冲区;Learner 进程持有两个缓冲区与优化器，负责混合采样、计算 SAC 损失、更新 actor/critic 与温度参数，并定期把最新权重同步给 Actor。两侧通过 gRPC 通讯，使环境采样与梯度更新完全解耦——Learner 以最快速度执行梯度步，Actor 只需按固定的物理控制频率稳定收集转移，梯度计算不会阻塞 env step。该架构同时为日后真机训练保留通路（虽本文真机阶段最终采用 BC 与 HG-DAgger，详见\cref{sec:real_task_deploy}）。

\begin{figure}[!htbp]
  \centering
  \includegraphics[width=0.7\textwidth]{fig_5_1_hilserl_arch.png}
  \caption{HIL-SERL 的 Actor-Learner 分布式架构。Actor 负责与环境交互、接收 HIL 干预、写入在线缓冲区;Learner 持有缓冲区与优化器，负责混合采样、SAC 更新并回推权重。两侧通过 gRPC 通讯，环境采样与梯度更新完全解耦。}
  \label{fig:hilserl_arch}
\end{figure}

\paragraph{仿真训练流程。}完整训练由两个在线阶段构成。\emph{Demo warmup 阶段}：训练开始时的前 500 步梯度更新仅从离线缓冲区采样（此时在线缓冲区尚空），用 50 条 demo 为两个 critic 与 actor 建立合理的初始值，避免在线交互起步时因价值函数随机初始化产生不稳定的策略行为。\emph{在线混合采样阶段}：warmup 结束后 Actor 进程开始与仿真环境交互并接收 HIL 干预。其中进入在线阶段后的前 100 个环境步被留作\emph{纯数据收集}（$\texttt{online\_step\_before\_learning}=100$，只向在线缓冲区写入转移而不触发梯度更新），其作用在于让在线缓冲区先积累一小批与当前策略分布匹配的样本，避免在缓冲区几乎为空时就以极度偏置的采样分布破坏 warmup 阶段已获得的初始值。之后进入稳态训练循环：每个环境步触发 2 次梯度更新（$\text{UTD}=2$），每次更新从离线与在线缓冲区各采 $|\mathcal{B}|/2$ 条转移;干预产生的转移按上一段所述回注离线缓冲区;每 4 个梯度步 Learner 把 Actor 与离散夹爪 critic 的最新权重通过 gRPC 同步推送给 Actor 进程。所有超参数见\cref{tab:hyperparams}（仿真 HIL-SERL 列）。

\section{备份策略仿真训练（SAC $+$ DR）}\label{sec:backup_sim_train}

本节给出备份避让策略的仿真训练全栈：问题定义、训练范式选择、场景设计、奖励函数演化、几何升级与域随机化。备份策略与\cref{sec:sim_task_train} 任务策略的最大区别是采用\emph{纯仿真 SAC $+$ DR} 而非 HIL-SERL——下文 \cref{sec:backup_paradigm} 给出范式选择的详细论证。

\subsection{问题定义与安全需求}\label{sec:backup_problem}

\textbf{备份避让策略的适用情境}：在机械臂执行抓取/放置任务的过程中，人手可能间歇性地进入工作空间——例如操作员调整物料位置、补充耗材、或在 HIL 训练期间手动干预。在这类间歇性人机共享情境下，备份策略\emph{在人手接近 TCP 时由其接管、主动避让、结束后交还任务策略}。该策略覆盖两种情境：(i)~生产部署时的人机协同安全（主要应用）;(ii)~\cref{ch:task_policy} HIL 训练期间操作者介入时的安全保障（派生应用）。

\paragraph{安全需求。}备份策略需同时满足四条要求：(1)~\textbf{避碰}——TCP 与人手距离 $> 8\,\text{cm}$;(2)~\textbf{保持夹持}——避让过程绝不放开已夹方块;(3)~\textbf{工作空间硬边界}——TCP 不越出预设笛卡尔工作空间（仿真 spawn box 与真机 \texttt{abs\_pose\_limit} 联合定义）;(4)~\textbf{原地闪避}——避让期间 TCP 累计位移 $\leq 15\,\text{cm}$，以保证任务策略能在原位附近恢复执行。违反任一条则 episode 即以 $-10$ 终止，具体奖励形式见\cref{sec:backup_reward}。

\paragraph{运行时策略切换机制。}Actor 侧同时持有任务策略与备份策略两套参数，由人手--TCP 欧氏距离决定当前步由哪一套产生动作：距离 $\geq 8\,\text{cm}$ 时任务策略动作生效，距离 $< 8\,\text{cm}$ 时立即切换到备份策略;备份策略至多连续执行 $10$ 步，期间若人手离开工作空间或步数用尽则切回任务策略。这一"距离触发 $+$ 时间有界"的硬切换把备份策略的训练目标简化为"在 $10$ 步内避让并稳定"，与任务层面的语义完全解耦。真机部署的具体阈值与几何对齐（Route A center-dist）见\cref{sec:backup_real_deploy}。

\subsection{训练范式选择}\label{sec:backup_paradigm}

\paragraph{设计演化：从硬编码冻结到分离式学习。}当前的分离式双策略设计经历过两次失败迭代，env 始终统一（任务策略与备份策略共享 \texttt{PandaPickPlaceSafeEnv}），变化的是如何"学习避让"。\emph{方案一}是\textbf{硬编码冻结}——人手接近时任务策略停止输出、机械臂保持构型。该方案静态安全但非学习型且被动低效（无法对运动模式差异化响应；人手沿原方向逼近时静止末端仍可能接触）。\emph{方案二}是\textbf{统一学习策略}——用同一个策略同时优化任务与避让，奖励沿用稀疏二值（成功 $= 1$、任何失败 $= 0$）。训练出现典型 reward hacking：策略学会人手出现时\emph{加速冲向目标板}抢在碰撞前完成放置，原因是该奖励不区分碰撞与超时（惩罚都是 0），"冲刺完成"在期望回报上优于"避让后完成"。

两次失败共同导出本节的最终架构：\textbf{env 不变（仍共享 25 维统一观测），把学习目标拆成两个专门策略}——任务策略保留原稀疏 $0/1$ 奖励（\cref{sec:task_reward}），仅优化"抓方块、放目标板"，对 $h_\text{active}$、$\vect{p}_\text{hand}$ 不产生避让响应;避让另开备份策略、采用非负即时 $+$ 大终止惩罚奖励（\cref{sec:backup_reward}）。两者由运行时距离判据硬切换——既保留硬编码冻结的安全下界，又通过学习型主动避让规避了"中断任务"与"被动接触"两种低效情形。

\paragraph{架构灵感：Kiemel 等人 2024。}本节的整体架构受 Kiemel 等人~\cite{kiemel2024safe} 启发：他们先用无模型 RL 训练一个 backup policy 产生避让动作，然后在学习其他任务时把它作为安全层——每步决策前估计碰撞风险，若风险过高则用 backup 的动作替换。本文沿用"预先训练 backup、运行时切换"的思路，并做两处调整：(i)~把使用场景从\emph{通用任务学习的安全层}换成 HIL 训练与人机共享工作空间下的协同安全;(ii)~把在线风险估计简化为基于人手--TCP 距离的硬切换，降低真机部署的计算与工程复杂度。奖励设计层面本文保留了其核心原则——\emph{非负即时奖励 $+$ 大终止惩罚}。

\paragraph{为何采用纯仿真 $+$ 域随机化。}与\cref{sec:sim_task_train} 任务策略的 HIL-SERL 形成对照，备份策略采用\emph{纯仿真 SAC $+$ 域随机化}而非 HIL 训练。四点理由：\emph{(i)~观测模态}——输入为低维状态、无 RGB 图像，不存在视觉 sim-to-real gap;低维观测的差距（位置噪声、速度抖动、检测延迟）由 DR 精确覆盖（\cref{sec:backup_training}）。\emph{(ii)~递归安全}——若备份本身也需 HIL 训练，则训练它时又需要一个更底层的"备份的备份"来保障操作者安全，无穷退化;纯仿真从根本上绕开这个问题。\emph{(iii)~可评估性}——仿真可以构造覆盖完整的障碍物运动场景，并以碰撞真值直接评估策略安全性。\emph{(iv)~偏差鲁棒性}——避让行为主要依赖障碍物与 TCP 的\emph{相对几何}而非绝对位置，DR 观测噪声使策略对大致方向而非精确位置作反应，因此 backup 对编码器偏差天然鲁棒、无需 bias-aware 训练;实证上，把无偏差环境下训练的备份 checkpoint 直接在带 joint 1 随机偏差的评估环境下运行，存活率与无偏差环境基本一致（详见\cref{sec:sim_experiments}）。

\subsection{场景、观测与运动模式}\label{sec:scenarios}

本节给出备份策略的仿真环境设计——两个训练场景、单帧观测向量的构成、帧堆叠策略，以及四种障碍物运动模式。

\paragraph{场景 S1 / S2 与观测向量。}本节在两个递进的仿真场景上训练备份策略：\emph{S1} 含一个间歇性出现的移动障碍物（模拟单只人手）;\emph{S2} 在 S1 基础上额外放置一个位置随机但固定不动的静止障碍物（模拟工作台上的固定物件）。两个场景共享相同的单帧观测结构——$18$ 维机器人状态（$7$ 关节角 $+$ $7$ 关节速度 $+$ $1$ 夹爪开度 $+$ $3$ TCP 位置）加每个障碍物的 $10$ 维观测（$1$ 维激活标志 $+$ $3$ 维绝对位置 $+$ $3$ 维差分速度 $+$ $3$ 维相对 TCP 位置）。S1 单帧 $28$ 维，S2 单帧 $38$ 维。激活标志刻画障碍物的间歇出现：障碍物在工作空间内为 $1$、离开后变为 $0$，从而策略可以在障碍物消失时回到正常状态、在下次出现时重新进入避让。

\paragraph{帧堆叠。}单帧观测已含障碍物的差分速度，但不同运动模式在时序上呈现出显著不同的特征——STOP\_GO 的周期性停顿意味着"当前速度 $\approx 0$"可能是真正停止也可能是短暂停顿;PASSING 的加速横穿与 LINEAR 的匀速直线只有在观察连续数步后才能可靠区分。为让策略基于一小段历史而非单一瞬时观测判断，本节对单帧观测做 $3$ 帧堆叠：
\begin{equation*}
  \vect{o}_t^{\text{stack}} = [\vect{o}_{t-2}^{\text{single}};\; \vect{o}_{t-1}^{\text{single}};\; \vect{o}_t^{\text{single}}],
\end{equation*}
从而 S1 的完整观测为 $28 \times 3 = 84$ 维，S2 为 $38 \times 3 = 114$ 维。这等价于在观测层面隐式提供了二阶运动信息（加速度与模式特征），使策略能识别运动模式并对短时未来运动作出预测。

\paragraph{障碍物运动模式与 episode 设定。}每个 episode 开始时从四种运动模式中均匀随机选择一种并在 episode 内保持不变（\cref{tab:motion_modes}），以覆盖人手接近的典型模式;策略须对全部四种模式保持鲁棒。训练 episode 长度固定为 $10$ 步（@ $10$\,Hz，即 $1$\,s），并采用\emph{无折扣}设定 $\gamma = 1.0$——在这样短的 episode 下每一步的安全性等价重要，任何折扣都会人为扭曲避让行为在时间维度上的权重。

\begin{table}[htbp]
  \centering
  \caption{备份策略训练中的四种障碍物运动模式}
  \label{tab:motion_modes}
  \footnotesize
  \begin{tabular}{@{}llll@{}}
    \toprule
    模式 & 轨迹 & 关键参数 & 模拟场景 \\
    \midrule
    LINEAR   & 匀速直线朝 TCP & --- & 人直接伸手取物 \\
    ARC      & 弧线绕向 TCP & 每步方向向量旋转 $\omega \sim \mathcal{U}(-0.1, 0.1)$\,rad & 人从侧面绕过 \\
    STOP\_GO & 朝 TCP 随机停顿 & 每步 $30\%$ 概率停顿 & 人犹豫后再伸手 \\
    PASSING  & 垂直于 TCP 方向横穿 & 速度 $\times\, 1.3$ & 人手路过而非针对机器人 \\
    \bottomrule
  \end{tabular}
\end{table}

\subsection{奖励函数设计演化}\label{sec:backup_reward}

备份策略的奖励函数经历三次迭代——V1 / V2 / V3，每次失败暴露出一种典型的 reward hacking 模式，共同导出 V3 的最终形式;V3 在真机部署阶段进一步演化为 V3b 几何（详见\cref{sec:v3b_geometry}）。本小节按时间顺序陈述三个版本并总结设计原则;具体训练曲线与量化对比属于实验部分，留至\cref{sec:sim_experiments}。

\paragraph{V1：接近度奖励及其失败。}V1 采用"远离障碍物"式的即时奖励——以 TCP 与最近障碍物的距离为驱动信号：
\begin{equation*}
  r_t^{\text{V1}} = r_{\text{prox}}(d_t^{\text{obs}}) - 0.01 \lVert \vect{a}_t \rVert,\quad
  r_{\text{prox}}(d) = \min\!\left(1, \, (d / d_{\text{danger}})^2\right),
\end{equation*}
其中 $d_t^{\text{obs}}$ 为 TCP 与最近障碍物之间的欧氏距离;碰撞时 $r_{\text{term}} = -1.0$。问题在于 $r_{\text{prox}}$ \emph{直接奖励远离}——策略学到的最优行为是尽可能跑得远，而不是留在原位小幅闪避。这与本节部署需求（\cref{sec:backup_problem} 安全要求 (4)"原地闪避"）直接冲突：策略"避让了人手"却"丢失了任务位置"。\emph{教训}：奖励必须与部署目标对齐——避让的目标不是"离障碍物越远越好"，而是"既不碰撞也不远离"。

\paragraph{V2：全负奖励及其失败。}V2 改为\emph{全负}软约束形式，以位移 $+$ 动作幅度 $+$ 动作平滑三项惩罚劝阻策略跑远。为书写紧凑引入两个缩写：$\Delta\vect{p}_t := \vect{p}_t^{\text{TCP}} - \vect{p}_0^{\text{TCP}}$（相对 episode 起点的累计位移），$\Delta\vect{a}_t := \vect{a}_t - \vect{a}_{t-1}$（动作增量），则每步奖励为
\begin{equation*}
  r_t^{\text{V2}} = -\, 0.5 \lVert\Delta\vect{p}_t\rVert - 0.01 \lVert\vect{a}_t\rVert - 0.01 \lVert\Delta\vect{a}_t\rVert,
\end{equation*}
碰撞终止 $r_{\text{term}} = -1.0$，$\gamma = 0.99$。问题在于\textbf{存活每一步都带来负奖励}——完整存活 10 步的总回报约为 $-0.5$ 到 $-0.8$ 之间，一旦第 1 步碰撞则为 $-1.0$，两者差距仅 $\sim 0.2$;位移积累稍大时存活总回报甚至比立即碰撞更糟。策略学到的最优行为变成了\emph{尽快终止}（"快死优于慢死"），训练日志中平均奖励由 $-0.88$ 持续恶化到 $-1.00$、episode 长度逐步缩短直至几乎每回合首步即碰撞（完整曲线见\cref{sec:sim_experiments}）。\emph{教训}：存活必须严格优于任何终止——即时奖励不能为负，且终止惩罚必须远大于任何可能的存活轨迹累积回报。

\paragraph{V3：非负即时奖励 $+$ 大终止惩罚。}V3 严格遵循 Kiemel 等人~\cite{kiemel2024safe} 的"每步即时奖励非负、终止给予大惩罚"原则，并在软奖励中保留位移项以实现"原地闪避"目标。存活每一步的即时奖励为
\begin{equation*}
  r_t^{\text{V3}} = 0.5 - 0.5 \lVert\Delta\vect{p}_t\rVert - 0.01 \lVert\vect{a}_t\rVert - 0.01 \lVert\Delta\vect{a}_t\rVert \;\geq\; 0,
\end{equation*}
在碰撞、方块掉落、TCP 越工作空间硬边界、或累计位移 $\lVert\Delta\vect{p}_t\rVert > 15$\,cm 任一触发时，episode 以大终止惩罚 $r_{\text{term}} = -10.0$ 结束;若 episode 安全存活满 $10$ 步则额外给予 $+5.0$ 完成奖励;折扣 $\gamma = 1.0$。\cref{tab:v3_returns} 列出几种典型场景的总回报：不动 $+10.0$ $>$ 小闪避 $\approx +9.7$ $>$ 大闪避 $\approx +8.5$ $\gg$ 碰撞或过度反应 $-10.0$。这一排序有三个关键性质：(i)~\emph{每步即时奖励非负}，存活永远优于提前终止，消除了 V2 的"快死"漏洞;(ii)~\emph{终止惩罚远大于任何存活轨迹的累积上限}（$|{-10.0}| \gg +10.0$），主动终止不是期望回报上的优解;(iii)~不动 $>$ 小闪避 $>$ 大闪避的软排序让策略在可以不动时倾向于不动（对齐"原地闪避"目标，消除了 V1 的"跑远"漏洞），不得不动时倾向于小幅动作。

\begin{table}[htbp]
  \centering
  \caption{V3 奖励下典型场景的总期望回报（$\gamma = 1.0$，$10$ 步 episode）}
  \label{tab:v3_returns}
  \footnotesize
  \begin{tabular}{@{}llc@{}}
    \toprule
    场景 & 行为描述 & 总回报 \\
    \midrule
    完美静止 & 原地不动、人手路过未接触 & $+10.0$ \\
    小幅闪避 & 累计位移 $\approx 5$\,cm 后存活 & $\approx +9.7$ \\
    大幅闪避 & 累计位移 $\approx 12$\,cm 后存活 & $\approx +8.5$ \\
    过度反应 & 累计位移 $> 15$\,cm 触发硬约束 & $-10.0$ \\
    碰撞失败 & 首步未闪避、与人手相撞 & $-10.0$ \\
    \bottomrule
  \end{tabular}
\end{table}

\paragraph{设计原则总结。}V1 $\to$ V2 $\to$ V3 的三次迭代归纳出四条在避让类短 episode 任务上普遍适用的奖励设计原则：\textbf{(1)}~即时奖励必须\emph{非负}，保证存活严格优于任何终止;\textbf{(2)}~终止惩罚必须\emph{远大于}任何存活轨迹的累积上限，防止"主动终止"成为最优解;\textbf{(3)}~位移等跨 episode 的硬约束必须用\emph{终止惩罚}而非软惩罚表达，避免通过累积软惩罚规避终止;\textbf{(4)}~短 episode 任务应使用\emph{无折扣}设定 $\gamma = 1.0$，避免把安全权重在时间维度上做不必要权衡。

\subsection{V3$\to$V3b 几何升级：从单球到 5 球全臂}\label{sec:v3b_geometry}

V3 奖励的几何模型只检测 \texttt{panda\_hand} 单球与障碍物的距离。真机部署阶段（\cref{sec:backup_real_deploy}）发现该模型存在一个失效模式：策略学会让 \texttt{panda\_hand} 避开障碍物，但\textbf{肘部与前臂仍可能撞向障碍物}——单球检测无法捕获 link3/link4 等中段链节的碰撞风险。本文据此把 V3 升级为 V3b 几何，包含三处改动：

\begin{enumerate}
  \item \emph{多球全臂检测}：单球（仅 \texttt{panda\_hand}，半径 $0.10$\,m）扩展为 $5$ 球（\texttt{link3, link4, link5, link6, hand}，半径分别为 $0.07, 0.07, 0.065, 0.065, 0.10$\,m）。碰撞判定取所有球的最小距离，使训练奖励直接覆盖肘/前臂触障风险。
  \item \emph{障碍物半径与 TCP-障碍物距离调整}：障碍物半径从 $0.035$\,m 提升至 $0.10$\,m（对齐真机 hand bbox 等效半径）;障碍物 spawn 距离从 $\mathcal{U}(0.21, 0.30)$\,m 提升至 $\mathcal{U}(0.30, 0.40)$\,m;手速从 $\mathcal{U}(0.005, 0.015)$\,m/step 提升至 $\mathcal{U}(0.015, 0.030)$\,m/step;最大允许位移从 $0.30$\,m 提升至 $0.50$\,m。
  \item \emph{Saturating proximity bonus}：在 V3 reward 上加入饱和接近奖励（饱和值 $0.20$ @ clearance $0.10$\,m），让策略对临界距离（接近但不碰撞）有正梯度信号，避免 V3 在多球检测下因"远离即可"训不出精细避让。
\end{enumerate}

V3b 配套的 FSM 阈值升级见\cref{sec:backup_real_deploy} （Route A center-dist：$D_\text{SAFE}$ 0.30 $\to$ 0.40\,m，$D_\text{CLEAR}$ 0.35 $\to$ 0.45\,m）。V3b 训完版（\texttt{backup\_policy\_s1\_v3b\_300k}）在真机首选 ckpt 中存活率 95.0\%，相对 V3 单球版 71.5\% 显著提升;\texttt{hand\_collision} 率从 17.5\% 降至 3.0\%（详见\cref{sec:sim_experiments}）。

\subsection{训练协议与域随机化}\label{sec:backup_training}

\paragraph{网络架构（相对\cref{sec:network} 的简化）。}备份策略的网络主体与\cref{sec:network} 完全同构——SAC 的连续 Actor 与 Twin Critic 都以隐层 $[256, 256]$、激活 ReLU 的 MLP 为主体，学习率、梯度裁剪、目标网络软更新等基础超参均复用任务策略的设定。动作空间为 $6$ 维连续末端增量 $(\Delta x, \Delta y, \Delta z, \Delta r_x, \Delta r_y, \Delta r_z)$、单步缩放 $0.03$\,m，夹爪在备份策略激活期间始终锁定、不参与动作。两处简化：(i)~输入为低维状态、\emph{无视觉编码器};(ii)~采用\emph{纯在线 SAC}，无离线 demo 缓冲区、不做 RLPD 式的 demo bootstrap。其余非零差异列于\cref{tab:backup_hparams}。其中 UTD 比例取 $8$（相对任务策略的 $2$ 显著提高），是对"短 episode、在线数据稀少"的一种补偿;训练过程中也试过 UTD $\in \{2, 4\}$ 等设定，对最终存活率基本无影响，仅影响训练收敛的 wall-clock 速度。

\begin{table}[htbp]
  \centering
  \caption{备份策略相对任务策略的关键超参数差异}
  \label{tab:backup_hparams}
  \footnotesize
  \begin{tabular}{@{}lll@{}}
    \toprule
    超参数 & 任务策略（\cref{ch:task_policy}） & 备份策略（本节） \\
    \midrule
    离线 demo 缓冲区 & $50$ episodes & 无 \\
    总在线训练步数 & $10^5$ & $2 \times 10^5$（V3b 训 $3 \times 10^5$） \\
    UTD 比例 & $2$ & $8$（V3b 调回 $4$ 防过拟合） \\
    SAC 初始温度 $\alpha_0$ & $0.01$ & $0.1$ \\
    在线缓冲区容量 & $10^5$ & $2 \times 10^5$ \\
    Episode 长度 & $200$ 步 & $10$ 步（V3b $20$ 步） \\
    折扣因子 $\gamma$ & $0.97$ & $1.0$ \\
    \bottomrule
  \end{tabular}
\end{table}

\paragraph{域随机化参数。}域随机化\emph{仅作用于观测层}，碰撞检测始终使用真实位置——这是安全边界的底线：DR 让策略对感知噪声鲁棒，但不应让策略在"假碰撞"上练习避让。具体参数见\cref{tab:backup_dr}，各项的真机来源将在下段 sim-to-real 对齐分析中一一对应。

\begin{table}[htbp]
  \centering
  \caption{备份策略训练中的域随机化参数}
  \label{tab:backup_dr}
  \footnotesize
  \begin{tabular}{@{}lll@{}}
    \toprule
    参数 & 分布 & 模拟内容 \\
    \midrule
    障碍物位置噪声 & $\mathcal{N}(\vect{0},\,0.03^2 \mathbf{I}_3)$\,m & 视觉 $+$ 深度的累计定位误差 \\
    障碍物速度噪声 & $\mathcal{N}(\vect{0},\,0.01^2 \mathbf{I}_3)$\,m/step & 帧间差分抖动 \\
    观测延迟 & $\mathcal{U}\{0, 1, 2\}$ 步 & 视觉检测 $+$ 推理延迟 \\
    障碍物生成距离 & $\mathcal{U}(0.30,\, 0.40)$\,m （V3b）& 不同的进入距离 \\
    障碍物移动速度 & $\mathcal{U}(0.015,\, 0.030)$\,m/step （V3b）& 不同的手速 \\
    TCP 位置噪声 & $\mathcal{N}(\vect{0},\,0.005^2 \mathbf{I}_3)$\,m & 外部 TCP 定位的单帧误差 \\
    \bottomrule
  \end{tabular}
\end{table}

\paragraph{Sim-to-Real 对齐分析。}\cref{tab:backup_sim2real} 按观测与控制链上的每一个环节比较仿真与真机间的差距。核心观察是：\textbf{机器人自身状态} $(\vect{q}, \dot{\vect{q}}, \text{TCP})$ 由编码器直接读出，在 sim 与 real 之间没有 gap（这也正是本节选择纯状态观测的主要动机）;而所有与\textbf{障碍物感知}相关的维度（位置、速度、延迟）都依赖外部视觉系统，必须由 DR 覆盖。控制器从仿真的 OSC 到真机的笛卡尔阻抗控制存在结构差异，但在小步 TCP 增量动作空间上近似等价，严格论证见\cref{ch:real_chapter}。

\begin{table}[htbp]
  \centering
  \caption{备份策略的 Sim-to-Real 差距分析}
  \label{tab:backup_sim2real}
  \footnotesize
  \begin{tabular}{@{}llll@{}}
    \toprule
    环节 & 仿真来源 & 真机来源 & Gap 弥补方式 \\
    \midrule
    机器人状态 $\vect{q}, \dot{\vect{q}}, \text{TCP}$ & MuJoCo 真值 & 关节编码器 & 无 gap \\
    手部位置 & MuJoCo 真值 & 视觉检测 $+$ D455 深度 & DR 位置噪声（详见\cref{sec:hand_detection_evolution}） \\
    手部速度 & 真值差分 & 视觉结果帧间差分 & DR 速度噪声 \\
    观测延迟 & $0$ & $\sim 10$\,ms 检测延迟 & DR 延迟 $\mathcal{U}\{0,1,2\}$ 步 \\
    坐标变换 & 无 & $T_\text{cam}\!\to\!\text{robot}$ 标定 & 含入位置噪声项 \\
    控制器 & OSC & 笛卡尔阻抗控制 & 小步增量下近似等价（\cref{ch:real_chapter}） \\
    障碍物运动模式 & 4 种预设模式随机 & 真实人手动作 & 4 种模式覆盖典型动作 \\
    \bottomrule
  \end{tabular}
\end{table}

手部检测器的具体选型（MediaPipe 初选 $\to$ WiLoR 部署演化）属于真机部署决策，搬到\cref{sec:hand_detection_evolution} 详述。

\section{仿真实验与分析}\label{sec:sim_experiments}

本节给出本章两条主线（任务策略 HIL-SERL $+$ 备份策略 SAC $+$ DR）对应的全部仿真实验数据。所有实验均在 PandaPickCube / PandaPickPlaceSafe / PandaBackupPolicy 系列 MuJoCo 环境中完成，每个偏差值评估 $50$ episodes，每组超参数训练 $3$ 个 seed。

\subsection{实验设置}\label{sec:sim_exp_setup}

\paragraph{环境。}基于 MuJoCo 3.0 的 7-DoF Franka Panda 仿真，OSC 控制 $10$\,Hz、物理仿真 $500$\,Hz。任务策略评估在 PandaPickCube 上完成：抓取桌面方块、抬升 $20$\,cm 即视为成功。备份策略评估在 PandaBackupPolicyS1 / S2 上完成：每 episode 等概率从 \cref{tab:motion_modes} 四种模式随机选一种障碍物运动，存活满 $10$ 步且未触发任何终止条件视为成功。

\paragraph{统计协议。}每组实验 $3$ seeds，每 seed 评估 $50$ episodes，报告 mean $\pm$ std;偏差扫描时每个偏差值独立 $50$ episodes。训练超参除特别说明外按\cref{tab:hyperparams} 与\cref{tab:backup_hparams} 默认设定。

\subsection{H1：编码器偏差退化基线}\label{sec:exp_h1}

\textbf{假设。}无偏差训练得到的策略在编码器偏差下性能退化。

\textbf{实验。}训练无偏差策略（PickCube 18D 观测，无 \texttt{real\_tcp}），在 $b \in \{0.0, 0.05, \ldots, 0.30\}$\,rad 下分别评估。评估环境采用\cref{alg:sim_step} 同时满足\cref{prop:injection} 条件 (C1)（控制律使用 $\vect{q}_\text{true}+\vect{b}$）与 (C2)（观测返回 biased 状态）的双注入仿真步，确保偏差注入忠实复现\cref{sec:causal_chain} 的因果链。

\textbf{结果}（\cref{tab:h1_degradation}）：偏差 $> 0.10$\,rad 起成功率显著下降，$\geq 0.25$\,rad 接近完全失败;末端位置随偏差线性偏移（每 $0.05$\,rad 约 $3$\,cm）。这一退化曲线直接对应\cref{prop:non_identifiable}：基础观测下偏差不可辨识，策略无法对不同 $\vect{b}$ 做差异化补偿。

\begin{table}[htbp]
  \centering
  \caption{H1：无偏差策略在不同 J1 偏差下的退化}
  \label{tab:h1_degradation}
  \footnotesize
  \begin{tabular}{@{}cccc@{}}
    \toprule
    Bias (rad) & 成功率 & 平均步数 & TCP 偏移 \\
    \midrule
    0.00 & 100\% & 13.5 & 0\,cm \\
    0.05 & 100\% & 13.5 & $\sim 3$\,cm \\
    0.10 & 80\%  & 55.9 & $\sim 6.5$\,cm \\
    0.15 & 62\%  & 62.5 & $\sim 10$\,cm \\
    0.20 & 51\%  & 99.8 & $\sim 13$\,cm \\
    0.25 & 5\%   & 184.5 & $\sim 16$\,cm \\
    0.30 & 0\%   & 191.1 & $\sim 19$\,cm \\
    \bottomrule
  \end{tabular}
\end{table}

\subsection{H2 $+$ H3：固定 vs 随机偏差训练}\label{sec:exp_h2}

\textbf{假设。}固定偏差训练导致单点条件化与零偏差过补偿;随机偏差训练实现全范围鲁棒。

\textbf{实验。}三组对照训练（PickCube 24D 观测，含 \texttt{real\_tcp}）：(a) 无偏差;(b) 固定 $b = 0.2$\,rad;(c) 随机 $b \sim \mathcal{U}[0, 0.25]$\,rad。各自训完后在全偏差扫描上评估。

\textbf{结果}（\cref{tab:h2h3_three_strategies}）：
\begin{itemize}
  \item \emph{无偏差策略}：$b = 0$ 完美但 $b > 0.10$ 起退化（同 H1）。
  \item \emph{固定偏差策略}：训练点 $b = 0.2$ 附近 $99\%$+，但 $b = 0$ 时降至 $83\%$——\textbf{零偏差过补偿}的直接实证，对应\cref{subsec:identifiability_vs_solvability} 的"单点条件化 $\to$ 过补偿"论断。
  \item \emph{随机偏差策略}：全范围 $92$–$100\%$ 鲁棒，且\textbf{超训练范围泛化}（$b = 0.30$\,rad 仍 $99\%$，训练上限 $0.25$）。
\end{itemize}

\begin{table}[htbp]
  \centering
  \caption{H2 $+$ H3：三种策略在偏差扫描下的成功率对比（PickCube 24D 环境;与\cref{tab:h1_degradation} H1 退化曲线同一 bias 网格）}
  \label{tab:h2h3_three_strategies}
  \footnotesize
  \begin{tabular}{@{}cccc@{}}
    \toprule
    Bias (rad) & 无偏差策略 & 固定 $b=0.2$ 策略 & 随机 $b\in[0,0.25]$ 策略 \\
    \midrule
    0.00 & 100\% & 83\%   & 92\% \\
    0.05 & 100\% & 87\%   & 100\% \\
    0.10 & 80\%  & 99\%   & 99\% \\
    0.15 & 62\%  & 100\%  & 96\% \\
    0.20 & 51\%  & 100\%  & 97\% \\
    0.25 & 5\%   & 100\%  & 96\% \\
    0.30 & 0\%   & 99\%   & 99\% \\
    \bottomrule
  \end{tabular}
\end{table}

\subsection{观测空间消融：从 18D 到 24D}\label{sec:exp_ablation}

\textbf{假设。}观测维度的可辨识性结论可被消融实验直接验证——剔除独立外部 TCP $\vect{o}_\text{real}$ 后偏差不可辨识、训练失败。

\textbf{实验。}三档观测空间对照（PickCube，随机偏差训练）：18D（无 \texttt{real\_tcp} 与 \texttt{block\_pos}）/ 21D（仅加 \texttt{block\_pos}）/ 24D（再加 \texttt{real\_tcp}）。

\textbf{结果}（\cref{tab:ablation_obs}）：
\begin{itemize}
  \item \emph{18D}：仅 robot proprio，全偏差扫描平均 $42\%$——基础观测下偏差不可辨识，对应\cref{prop:non_identifiable}。
  \item \emph{21D}：加目标位置但仍无独立 TCP，训练不收敛——"知目标不知己"，无法补偿编码器偏差。
  \item \emph{24D}：再加 \texttt{real\_tcp} 后达到 $92$–$100\%$——\cref{prop:identifiable} 的可辨识性条件被满足，策略隐式恢复 $\vect{b}$ 并条件补偿。
\end{itemize}

\begin{table}[htbp]
  \centering
  \caption{观测空间消融：理论预测与实测的一一对照}
  \label{tab:ablation_obs}
  \footnotesize
  \begin{tabular}{@{}lllll@{}}
    \toprule
    观测 & 维度 & 缺少的关键信号 & 结果 & 理论预测 \\
    \midrule
    18D & 18 & \texttt{real\_tcp} $+$ \texttt{block\_pos} & $42\%$ avg & 不可辨识（\cref{prop:non_identifiable}） \\
    21D & 21 & \texttt{real\_tcp} & 训练失败 & 知目标不知己 \\
    24D & 24 & （完整） & $92$–$100\%$ & 近似可辨识（\cref{prop:identifiable}） \\
    \bottomrule
  \end{tabular}
\end{table}

\paragraph{说明。}本组消融在 PickCube 24D 上做以隔离编码器偏差因素;\cref{sec:obs_design} 描述的部署 25D 观测是 24D $+$ hand 通道的等价扩展（hand 通道不参与可辨识性论证，仅服务运行时 backup 切换），消融结论可类比扩展到 25D 部署观测。

\subsection{Backup S1 V3b 主存活率}\label{sec:exp_backup_s1}

\textbf{实验。}在 PandaBackupPolicyS1 V3b 几何（\cref{sec:v3b_geometry}）上训练 $3 \times 10^5$ 步，max\_disp $= 0.50$\,m;在 $200$ episodes（每模式 $50$ episodes）上评估。

\textbf{结果。}V3b 训完版（\texttt{backup\_policy\_s1\_v3b\_300k}）：
\begin{itemize}
  \item \textbf{满存活率}：$95.0\%$（vs V3 单球版 $71.5\%$，vs V2 单球版 $100\%$ 但真机存在肘撞）
  \item \textbf{失败模式分布}：\texttt{hand\_collision} $3.0\%$（V3 单球版 $17.5\%$，下降 $5\times$）;\texttt{excessive\_displacement} $5.5\%$
  \item \textbf{运动模式分项}：LINEAR / ARC / STOP\_GO / PASSING 全部 $\geq 90\%$，最弱 ARC $90.5\%$（残余 $9.5\%$ 失败主要集中于此模式，机理待\cref{sec:sim_to_real_bridge} 后的真机评估实证）
\end{itemize}

V3b 几何升级\textbf{显著缓解}（hand\_collision $5\times$ 下降）但\emph{未完全消除}真机部署观察到的肘/前臂撞击失效模式（\cref{sec:v3b_geometry} 已展开）;V3 单球版作为对照保留，证明几何升级的方向是必要的，但 ARC 模式残余 $\sim 9.5\%$ 失败的细分（究竟是哪个 link 球撞击）留作 future work。

\subsection{Backup 偏差鲁棒性 zero-shot}\label{sec:exp_backup_bias}

\textbf{假设。}备份策略对编码器偏差天然鲁棒、无需 bias-aware 训练（\cref{sec:backup_paradigm} 第四点理由的实证）。

\textbf{实验。}用 \cref{sec:exp_backup_s1} 的 S1 baseline checkpoint（无偏差训练），分别在两个评估环境跑 $200$ episodes：(a) PandaBackupPolicyS1（无偏差）;(b) PandaBackupPolicyS1BiasJ1（J1 随机偏差 $b_1 \sim \mathcal{U}[0, 0.25]$\,rad）。

\textbf{结果。}存活率 baseline $99.0\%$ vs BiasJ1 $98.5\%$，差距 $0.5\%$ 在 $200$ ep 抽样噪声范围内（$\pm 1.5\%$ at $95\%$ CI）。

\textbf{解读限定（sanity check 而非鲁棒性结果）。}本实验作为\emph{解耦验证（sanity check）}而非外加鲁棒性证明：备份策略只读 \texttt{tcp\_true}（外部独立通道），编码器偏差对其观测层物理不可见，故"零退化"是设计预期而非习得的偏差鲁棒性。该结果实证的是"backup 观测通路与编码器路径已解耦"，因此可省去 bias-aware 训练成本;真正的偏差鲁棒性证明仍在\cref{sec:exp_h2}（H2/H3 任务策略实验）。

\subsection{Reward V1 / V2 / V3 演化训练曲线}\label{sec:exp_reward_evolution}

\textbf{实验。}三组训练在同一 PandaBackupPolicyS1 环境、同一 SAC 超参下使用三种奖励：V1（接近度）/ V2（全负）/ V3（非负 $+$ 大终止）。每组 $3$ seeds、$2 \times 10^5$ 步。

\textbf{记录指标。}V1：平均 TCP 累计位移（应逼近 $15$\,cm 上限）;V2：平均 episode 奖励（应从 $\sim -0.88$ 持续恶化到 $\sim -1.00$）+ 平均 episode 长度（应缩短）;V3：平均 episode 奖励（应稳定收敛到 $\sim +9$ 到 $+10$）+ 存活率（应升至 $\sim 99\%$）。

\textbf{结果。}三组训练曲线印证 \cref{sec:backup_reward} 的设计原则：
\begin{itemize}
  \item \emph{V1}：策略学会跑远，平均位移逼近 $15$\,cm 上限，触发 \texttt{excessive\_displacement} 大量终止——典型的 reward hacking。
  \item \emph{V2}：平均奖励 $-0.88 \to -1.00$ 单调恶化，episode 长度缩短到几乎首步即碰撞——"快死优于慢死"。
  \item \emph{V3}：稳定收敛到 $+9$ 到 $+10$ 区间，存活率升至 $\sim 99\%$，对应\cref{sec:exp_backup_s1} 的最终结果。
\end{itemize}

\subsection{Backup S2 多障碍场景扩展}\label{sec:exp_backup_s2}

\textbf{实验。}PandaBackupPolicyS2（移动 $+$ 静止双障碍），$114$D 观测（$38 \times 3$ 帧堆叠），UTD $= 4$（UTD $= 8$ 训 $1 \times 10^5$ 步崩溃），$2 \times 10^5$ 步。

\textbf{结果。}存活率 $83$–$91\%$（两次 eval 方差较大）;\texttt{hand\_collision} 率 $13.5\%$（vs S1 baseline $0.5\%$，升 $26\times$）;最弱模式 ARC $75.0\%$ / STOP\_GO $81.8\%$、最强模式 PASSING $91\%$。从单障碍 $99\%$ 到多障碍 $\sim 85\%$ 是\textbf{多障碍几何本身更难}（逃生锥变窄、样本组合从 $50$ 升到 $2500$），不是策略学不好。证明本节架构在多障碍场景上仍可扩展。

\section{仿真到真机的桥接}\label{sec:sim_to_real_bridge}

上一节给出的仿真闭环为下一章真机部署提供两点凭据，但需对其作用范围作严格限定。\textbf{可辨识性命题（\cref{prop:identifiable}）已在\cref{sec:exp_ablation} 24D 消融中实证为信息论上"成功的信号路径"}——但需明确：真机首跑只取该路径中的\emph{特权一半}（外部独立 TCP \texttt{tcp\_true}），暂未投入被污染的编码器路径，故\cref{ch:real_chapter} 的成功率矩阵验证的是"sim 已证可行的成功信号路径在真机上仍能跑通"，而\emph{并非}\cref{prop:identifiable} "biased $\vect{q}+\vect{b}$ 与 $\vect{o}_\text{real}$ 联合可辨识"的真机闭环;后者的两步消融（先验证 \texttt{tcp\_biased} 单独失败、再验证 $+\vect{q}+\vect{b}$ 恢复成功）留作\cref{ch:conclusion} 的 future work。换言之，\textbf{仿真侧 H2/H3 已闭环；真机侧仅取该路径的"clean half"，完整 identifiability path 在真机的闭环未完成}。

\textbf{Backup 策略 \texttt{backup\_policy\_s1\_v3b\_300k} 直接用作真机 ckpt}（不重训），配套\cref{sec:backup_real_deploy} 的 Route A center-dist 几何对齐与 V3b 阈值 $D_\text{SAFE} = 0.40$\,m。需注意 V3$\to$V3b 几何升级是真机部署观察到的失效模式（V3 单球版肘/前臂撞）倒逼仿真升级，而真机 V3b 评估见\cref{subsec:exp_backup_real}\emph{尚为待数据状态}——sim V3b $95\%$ 存活率构成 real 评估的预期上界，gap 大小待测。HIL-SERL 在仿真中稳定收敛但真机失败的对照分析见\cref{subsec:hilserl_failure}（真机 negative result 复盘）。

下章\cref{ch:real_chapter} 给出真机部署的完整栈：硬件架构与 B$+$D 偏差注入、任务策略真机部署（BC $+$ HG-DAgger）、备份策略真机部署（WiLoR $+$ Route A FSM）、以及三任务联合 fault injection 4 列成功率矩阵——本论文的核心实验数据。
